\chapter{原子核也会改变}

\begin{kaichang}
很多人家的天花板角落都挂着一只白色的小圆盘——烟雾报警器。把它翻过来，你可能会在铭牌上看到一行小字：“内含微量放射性物质镅-241”。放射性！就是那种让人联想到核辐射标志、防护服和警报声的东西，居然正大光明地住在你家卧室门口？

再看一眼你手里的香蕉。香蕉富含钾，而天然钾里混着极少量的钾-40——它也是放射性的。也就是说，此刻你的手里、你的身体里，正有原子核在悄悄地“变身”。

先别慌，也别猜。这一章结束时，你会知道烟雾报警器为什么必须装一颗会衰变的原子核，也会知道那只香蕉到底该不该让你担心。
\end{kaichang}

\section{一种看不见的“泄漏”}

1896 年的巴黎，物理学家贝克勒尔把几块含铀的矿石和一包用黑纸裹得严严实实的照相底片一起锁进了抽屉。他本来想做的是荧光实验，那几天气不好，实验搁置了。几天后他顺手把底片冲洗出来，却发现底片上印着矿石清晰的轮廓——底片被“曝光”了，尽管它没见过一丝光。

这是人类第一次撞见一种完全看不见、听不见、闻不到的“泄漏”：铀矿石在不停地向外发射某种东西，它能穿透黑纸，能让底片感光。后来居里夫妇发现，含铀矿渣的“泄漏”比纯铀还强，于是他们从一吨吨矿渣里提炼出两种新元素——钋和镭，并给这种现象起了名字：放射性。

今天，你不用碰任何矿石也能“听见”它。把一台盖革计数器（一种专门探测射线的仪器）靠近含镭的旧表盘或者某类矿物标本，扬声器会发出一连串咔哒声，像急促的雨点打在铁皮上。每一声咔哒，都是一个原子核在那一瞬间“变身”时被仪器逮个正着。

这里必须把安全话说在前面，而且只说一遍，但每一句都算数：\textbf{本章提到的所有放射性内容，只许阅读、观看视频资料或看老师演示，绝对禁止在家尝试。}不要拆烟雾报警器，不要收藏来路不明的“夜光”旧物件或矿物，不要网购任何标注“放射源”的东西。射线看不见、闻不到，你没有任何天生的报警器——这正是它的危险所在，也是它必须交给专业设备和专业人员的原因。

好消息是，你马上就要学到一套比报警器更好用的东西：一套判断“什么情况下真的危险、什么情况下纯属吓唬自己”的思维工具。

\section{原子核里的拔河比赛}

为什么有的原子核会“泄漏”，有的却几万亿年纹丝不动？答案藏在原子核内部一场永不停歇的拔河比赛里。

回忆第 6 章的图景：原子核只有原子的十万分之一大，却塞着原子几乎全部的质量。里面是两种粒子——带正电的质子，和不带电的中子。麻烦立刻就来了：同种电荷互相排斥，一堆质子挤在一个那么小的空间里，电磁排斥力拼命想把它们推开。可原子核没有散架，因为当粒子靠得足够近时，另一种强大得多的力开始起作用——强相互作用，它像超级胶水一样把质子和中子粘在一起。

关键在于，这两股力的“脾气”完全不同：

\begin{itemize}[itemsep=2pt]
\item 电磁排斥力作用范围远，核里每一个质子都在排斥其他所有质子，核越大，排斥越积越多。
\item 强相互作用只在紧挨着的粒子之间起作用，范围极短，出了“一臂之遥”就管不着了。
\end{itemize}

于是原子核的命运就成了两股力的比分。小的核里，邻居少而亲密，胶水占上风，核很稳定。核越大，电磁排斥的“总分”涨得越快——因为它在全体质子之间累加，而胶水只在邻居之间生效。超过某个临界点，再强的胶水也压不住场面，核就变得不稳定：它可能某一天突然吐出一小块碎片，或者调整一下内部的人员构成，变成一个能量更低、更安稳的新核。这就是放射性衰变——\textbf{不稳定原子核自发地改变自身结构、同时放出射线的过程}。

注意一个和第 7 章衔接的细节：同一种元素（质子数相同）可以有好几种同位素，差别在中子数。中子也参与“胶水”却不受电磁排斥，所以中子数是否“配比合适”，直接决定这个核稳不稳。碳-12（6 个质子、6 个中子）稳如泰山，碳-14（6 个质子、8 个中子）却是放射性的——同一张元素身份证，两种截然不同的核命运。

\section{三种衰变，各改各的账}

不稳定的核“变身”主要有三条路径，对应三种射线。它们改变的账目完全不同，值得一条条算清楚。

\textbf{$\alpha$ 衰变：吐出一个氦核。}太大的核常常一次抛出 2 个质子和 2 个中子捆成的小包——那恰好是一个氦原子核。抛出后，母核的质量数减 4，质子数减 2。质子数就是元素身份证（第 7 章），所以\textbf{元素变了}。烟雾报警器里的镅-241 走的就是这条路：
\[{}^{241}_{95}\chem{Am} \rightarrow {}^{237}_{93}\chem{Np} + {}^{4}_{2}\chem{He}\]
镅变成了镎。$\alpha$ 粒子个头大、电荷多，冲不出多远——一张纸、几厘米空气、皮肤的表层就能挡住它；但如果放出它的物质被吃进或吸进体内，近距离的破坏就很严重。

\textbf{$\beta$ 衰变：中子改行当质子。}中子偏多的核里，一个中子会转变成一个质子，同时放出一个高速电子（$\beta$ 粒子）和一颗几乎探测不到的反中微子。质量数不变，质子数加 1——\textbf{元素又变了}，往后挪了一格。考古学家依赖的碳-14 就是这样变成氮的：
\[{}^{14}_{6}\chem{C} \rightarrow {}^{14}_{7}\chem{N} + {}^{0}_{-1}\chem{e}\]
$\beta$ 粒子比 $\alpha$ 小得多，能穿透皮肤表层，但几毫米铝板就能拦住。

\textbf{$\gamma$ 衰变：只放能量，不改成分。}上面两种变身之后，新核往往还处在“激动”的高能状态，它会放出一束能量极高的电磁波——$\gamma$ 射线——平静下来。质子和中子数都不动，元素不变。但 $\gamma$ 射线穿透力最强，需要厚厚的铅板或混凝土才能明显削弱，医学和工业上用的射线多半是它。

把三条路径并排看，你会发现一件足以让古人瞠目的事：\textbf{原子核是会改变的，元素是会互相转变的}。中世纪炼金术士梦想把铅变成金，穷尽一生也没能动摇原子核分毫——因为化学反应只动核外的电子，原子核在化学变化中安然无恙。而在衰变里，嬗变真实地在发生：铀经过一连串衰变，最终真的会变成稳定的铅。只是方向不由人挑选，速度也不由人控制。顺便记下一笔：恒星内部用另一种方式（核聚变）把轻元素真正“炼”成了重元素，你身上的碳和铁就是这样来的——这是第 16 章的故事。

\section{半衰期：不讲个体，只讲群体}

下一个问题很自然：一个碳-14 原子核，什么时候衰变？

答案是：\textbf{没有人知道，也没有任何办法知道}。不是因为仪器不够好，而是核衰变在本质上就是随机的——明天衰变和一百年后衰变的同一个核，此刻没有任何区别。你盯着它看，不会催它快一点；把它冷冻、加热、通电、加压，统统无效。

但大自然在这里留了一扇窗：\textbf{单个不可预言，群体却精确得惊人}。一大堆同样的放射性核，每经过一段固定的时间，剩下的数目恰好减半。这段时间叫半衰期。碳-14 的半衰期约 5730 年：一克活的碳-14，5730 年后只剩半克，11460 年后剩四分之一，再往后八分之一、十六分之一……而碘-131 的半衰期只有 8 天，钋-214 只有 0.00016 秒，铀-238 则长达 45 亿年——比地球的年龄短不了多少。

\begin{qiaoliang}
为什么“每过相同时间恰好减半”？因为每个核在单位时间里衰变的概率是固定不变的。这和你抛硬币一样：每一枚硬币正面朝上的概率都是二分之一，于是抛一次，一大把硬币里大约一半“出局”；把剩下的再抛一次，又一半出局。核衰变相当于每个核每时每刻都在抛一枚这样的硬币。设原来有 $N_0$ 个核，经过 $t$ 时间、半衰期为 $T$，剩下的数目就是
\[N = N_0 \times \left(\frac{1}{2}\right)^{t/T}\]
拿碘-131（$T=8$ 天）算一笔：医院里一批碘-131 存放 80 天，$t/T=10$，剩下 $(1/2)^{10}=1/1024$，不到千分之一。这正是医院为什么敢用它治病——几周后病人体内几乎什么都不剩。
\end{qiaoliang}

做个真实的数量级估算，感受一下“群体”有多大。活生物体内的碳与环境保持交换，每克碳中约含五百亿个碳-14 原子。碳-14 半衰期 5730 年，约合 $1.8\times 10^{11}$ 秒，那么每秒衰变的比例约为 $\ln 2 \div (1.8\times 10^{11}) \approx 4\times 10^{-12}$。五百亿乘上这个比例：你身体里每克碳，每分钟大约有 14 个碳-14 核衰变。换句话说，\textbf{你本人就是一台微微作响的盖革计数器}——只是响动太小，完全无害。

\begin{shiyan}
\textbf{用一把硬币模拟半衰期}（安全，可在桌上完成）

材料：50～100 枚相同的硬币（或纽扣、围棋子），一个盒子，一张纸。

步骤：把硬币全部放进盒子，摇一摇倒出来。把所有正面朝上的硬币挑出来——它们“衰变”了，记录本批衰变数和剩余数。把剩余的硬币收回盒子，重复上述过程，直到剩个位数。每一轮代表“一个半衰期”。

观察点：（1）每轮大约淘汰剩余的一半，但绝不会精确到一半——请画出“剩余数—轮数”折线，看它是否呈一路减半的滑梯形。（2）如果只用 4 枚硬币重来，曲线还那么平滑吗？体会“规律只在大群体中精确”。（3）挑出任意一枚硬币：你能预言它第几轮出局吗？

安全提示：本实验没有任何放射性。真正的放射性观察（如盖革计数器、云雾室中粒子径迹）请只通过老师演示或正规科普视频进行。
\end{shiyan}

\begin{zhengju}
“元素会自己变成别的元素”——1902 年，当卢瑟福和他的助手索迪提出这个解释时，几乎等于当众撕掉化学的基石。当时的共识是：原子不可分、元素不可变，这是道尔顿以来化学的立身之本。

线索来自一串古怪的观察。居里夫妇发现镭之后，人们注意到放射性物质周围总会冒出一种具有放射性的气体（后来命名为氡），而氡自己又会慢慢“死掉”，同时在容器壁上留下一层新的放射性薄膜。卢瑟福和索迪把这串“放射性生放射性”的链条逐条追踪、逐个测量活性的消长，最后得出结论：这不是什么物质被“传染”了活性，而是原子核在一步步变身，每一步都变成一种新元素，同时放出射线。铀变钍、钍变镤、镤再变铀的另一种同位素……一路衰变到铅才停。

反对声很激烈，因为这套说法同时触犯了两条信条。但检验它的办法也干净利落：如果衰变真的造出新元素，那么铀矿里衰变的终点——铅——应该越老越多；氦（$\alpha$ 粒子抓住电子后的产物）也该在铀矿里大量积存。地质学家检测古老铀矿石，果然找到了与矿龄相称的铅和大量氦气泡。预言兑现，“嬗变”从异端变成了教科书。这个故事的要点不在谁聪明，而在于：哪怕结论听起来离经叛道，只要它给出能被检验的具体预言，科学就有办法裁决。
\end{zhengju}

\section{能量从哪来：裂变、聚变与结合能}

衰变放出的射线带着不小的能量，但核里还锁着一笔大得多的财富。要理解它，只需记住一句话：\textbf{把原子核拆散需要输入能量，反过来，把分散的粒子拼成一个核就会放出能量}——拼得越“划算”，放出的越多。

每种原子核的“划算程度”不同。把每个核子的结合能（平均每个质子或中子被粘得多牢）画成曲线，会看到一座中间高、两头低的山：山顶附近站着铁——它是所有核里粘得最牢的，能量上的“谷底”。这就指明了两个放能的方向：

\begin{itemize}[itemsep=2pt]
\item \textbf{核裂变}：比铁重得多的核（如铀-235）裂成两块中等大小的碎片，等于从山坡低处往谷底方向挪，每个核子都粘得更牢，差额以能量形式放出。一颗中子撞上铀-235 可触发裂变，裂变又抛出 2～3 颗新中子，新中子再去撞别的核——链条自己接下去，这就是链式反应。核电站的工作不是“烧核燃料”，而是用控制棒精确地拿捏这条链的速度，让能量以发热的形式匀速流出，再去烧水、推汽轮机发电。
\item \textbf{核聚变}：比铁轻得多的核（如氢的同位素氘、氚）合并成较重的核，同样往谷底挪，而且坡度更陡、放能更猛。太阳不是一堆火，而是一台持续运转了 46 亿年的聚变炉：每秒约有 6 亿吨氢聚变成氦。地球上的聚变发电至今仍是“永远还差几十年”的工程难题——把核加热到上亿度并关得住，比点火本身难得多。
\end{itemize}

为什么化学反应望尘莫及？因为化学只挪动核外电子，涉及的能量以“每个原子几个电子伏特”计；而核变化动的是核本身，涉及的能量以“每个核几百万电子伏特”计——相差约百万倍。这就是为什么一颗方糖大小的核燃料顶得上一卡车煤，也是为什么核武器的能量密度如此骇人。这同时也再次提醒我们：原子核是化学变化撼不动的禁区，本书后面所有章节讨论化学键、反应、代谢时，默认核都老老实实待在原地。

\section{放射性、剂量与风险：三件不同的事}

谈“核”之前，先把三个常被搅成一团的概念分开。分清了它们，你就不会再被“香蕉有放射性”这类标题吓到，也不会对真正的危险掉以轻心。

\textbf{放射性}是物质的性质：某种东西里，原子核每秒发生多少次衰变，单位是贝克勒尔（Bq），1 Bq 就是每秒一次衰变。一克镭约 $3.7\times 10^{10}$ Bq，一根香蕉约 15 Bq（来自钾-40），你自己的身体约几千 Bq（同样主要来自钾-40 和碳-14）。放射性描述的是“源”，还没说任何关于你的事。

\textbf{辐射剂量}是你实际吸收了多少能量、受到了多大影响。同样的源，隔着一堵铅墙和贴在皮肤上是两个世界；$\alpha$ 源拿在手里（射线穿不透皮肤表层）和吸进肺里也是两个世界。常用单位是希沃特（Sv），它把射线种类和器官敏感度都折算进去。给几个参照：天然本底（土壤、宇宙线、食物、氡气）让每人每年平均接受约 2～3 mSv；一次跨国飞行约 0.1 mSv；一次胸部 X 光约 0.02～0.1 mSv；吃一根香蕉约 0.0001 mSv。

\textbf{风险}是剂量换算成的健康后果的概率，而且关系并不简单：剂量很低时，风险小到淹没在日常生活风险里，无法分辨；剂量很高时（比如短时间内超过 1000 mSv），伤害确定且严重。中间地带靠大量人群统计来估计，充满不确定性。所以“有放射性”不等于“有危险”，正如“有汽车”不等于“会出车祸”——关键永远是多少剂量、以什么方式到达身体。医学上还有一条普遍原则：正当性和最优化——任何放射性检查都要收益明确、剂量能低则低。

分清这三件事之后，放射性的正当用途就不再神秘：

\begin{itemize}[itemsep=2pt]
\item \textbf{医学成像}：把微量短寿命放射性核素（如锝-99m，半衰期仅 6 小时）送进体内，它会聚集在特定器官，放出的 $\gamma$ 射线穿出身体被相机接收，医生看到的是一张“器官的功能照片”。PET 则利用正电子湮灭产生的 $\gamma$ 射线对，能看清肿瘤的代谢活跃程度。
\item \textbf{放射治疗}：癌细胞分裂旺盛，对辐射损伤的修复能力比正常细胞差。用精确聚焦的 $\gamma$ 射线束从多个角度照射，使剂量在肿瘤处叠加、在健康组织处分散——本质是把“剂量”这把刀磨得又快又准。
\item \textbf{测年}：活体不断与环境交换碳，体内碳-14 比例恒定；死亡后交换停止，碳-14 只减不增。测出残骸里碳-14 还剩几成，对照半衰期 5730 年，就能读出年代。这个方法把考古从“猜朝代”带进了“算年份”，适用范围大约五万年。
\item \textbf{核能}：裂变电站不排放二氧化碳，燃料能量密度极高；代价是必须管理长寿命的放射性废物和极低概率、极高后果的事故风险。它是工程问题，也是社会如何权衡“风险”的典型案例。
\end{itemize}

\section{这个模型的边界}

本章的图景很好用，但有三处边界值得说清。

第一，半衰期是统计规律，\textbf{对少量原子会失效}。说“碘-131 半衰期 8 天”，隐含的前提是“原子多得数不清”。如果你的硬币实验只剩 4 枚，下一轮可能一枚也不淘汰，也可能全军覆没——平均值依旧，单次结果无法预言。说“这种物质 80 天后完全消失”也是错的：理论上它永远还剩一个尾巴，只是少到无关紧要。

第二，“稳定”与“放射性”之间没有一刀切的线。有些曾被认为是绝对稳定的核，后来被发现其实也在衰变，只是半衰期长到 $10^{20}$ 年量级，比宇宙年龄还长几亿倍。稳定是个程度问题。

第三，本章把核当成“质子加中子加两股力”的粗粒模型，够用来理解衰变和能量，但解释不了全部细节——比如为什么某些质子数、中子数（所谓“幻数”）的核格外稳定，为什么核有形状、会振动。原子核内部的多体问题至今仍是活跃的研究前沿，没有完全解决。

\begin{wuqu}
“被辐射照过的东西，自己就会变成放射性的。”——这是流传最广的核误解之一，它让病人不敢做放疗，让人拒绝辐照灭菌的食品。事实是：$\gamma$ 射线、X 射线照过物体，能量被吸收，过程就结束了，不会在被照物体里留下任何“余孽”。医院拍过 X 光的你不带放射性，接受放疗的病人不带放射性，辐照过的香料和医疗器械也不带放射性。让普通物质“带上放射性”只有一种途径：它沾上了放射性物质的粉尘或液体（污染），或者用强中子流轰击它、把稳定核改造成放射性核（活化）——后者需要核反应堆级别的条件，日常生活中根本遇不到。要分清的是“照射”与“污染”：切过放射性样品的刀可能沾上放射性粉尘，需要按规范清洗；但站在源旁边被照了一下的刀，干干净净。
\end{wuqu}

\section{回到烟雾报警器和香蕉}

现在兑现开场的承诺。电离式烟雾报警器里那一点点镅-241（通常不到一微克），持续放出 $\alpha$ 粒子，把报警器内一个小气室里的空气分子打成带电离子，于是在两块电极之间形成一股稳定的微弱电流。烟雾颗粒一旦飘进气室，会吸附这些离子，电流立刻下降——电路检测到变化，警报响起。它“看见”烟，靠的是 $\alpha$ 衰变一刻不停的稳定输出。为什么选镅-241？因为它的半衰期长达 432 年，几十年内输出几乎不衰减，报警器才能常年免维护；又因为 $\alpha$ 射线穿不透塑料外壳，坐在它下面睡觉吸收的额外剂量，比一年多吃几根香蕉还小。当然，这不构成你拆开它的理由——内部的源片裸露时，风险完全是另一回事，请交给专门的回收渠道。

香蕉则是反面教材式的安慰剂：钾-40 半衰期约 12.5 亿年，一根香蕉每秒约 15 次衰变，剂量约 0.0001 mSv。你的身体含有约 140 克钾，其中钾-40 让你自带几千 Bq 的放射性——从你出生那天起就是如此，也必将伴随你一生。放射性不是现代工业发明出来攻击人类的怪物，它是这颗行星和这具身体的出厂配置。该敬畏的不是“放射性”三个字，而是剂量。

\section{核的故事还没有讲完}

这一章我们承认了一件不安的事实：原子核会变，元素会嬗变，连“稳定”都只是程度问题。但也正是这个不安的事实，给了人类癌症治疗的利器、深坑中文物的时钟、不冒黑烟的电站，以及理解宇宙的钥匙。

悬而未决的是那个更大的问题：如果原子核既能分裂也能合并、既在衰变也在生成，那么今天周期表上这一百来种元素，当初是从哪儿来的？大爆炸只造出了最轻的几种；你骨头里的钙、血液里的铁、手机里的稀土，全部是在恒星内部的聚变炉和恒星死亡时的剧烈爆发里，一次次核变化攒出来的。换句话说，本章讲的核物理，其实是宇宙制造元素的流水线。这条流水线怎样开工、怎样停工、又怎样把产品撒向太空——第 16 章《元素从哪里来，又去了哪里》会沿着核变化的线索一路追到恒星的生命尽头。

\begin{tiaozhan}
裂变一次放多少能量？算给你看。一次铀-235 裂变约放出 200 MeV（兆电子伏特），换算成焦耳：$200\times 10^{6}\times 1.6\times 10^{-19}\approx 3.2\times 10^{-11}$ J。看起来微不足道，但 1 克铀-235 约有 $6.02\times 10^{23}\div 235\approx 2.6\times 10^{21}$ 个原子，全部裂变放出的能量约 $8\times 10^{10}$ J——也就是 1 千克约 $8\times 10^{13}$ J。优质煤每千克燃烧约放出 $3\times 10^{7}$ J，两者相除：1 千克铀-235 完全裂变，相当于约 2700 吨煤。这个百万倍量级的差距，源头就是爱因斯坦那个著名的公式 $E=mc^{2}$：裂变后碎片总质量比原来的核轻了约千分之一，那一点点“消失”的质量按 $c^2$（一个天文数字）折算成了能量。化学反应中质量也变化，但只变十亿分之一量级，秤根本称不出来。
\end{tiaozhan}

\section*{练一练}

\begin{sikaoti}
\item 镭-226 经一次 $\alpha$ 衰变变成氡。写出这个核反应式（镭是 88 号元素，质量数减 4；查周期表确认产物是几号元素），并说明：产物的化学性质为什么会和镭完全不同？
\item 一种放射性核素的半衰期是 5 天。一批样品 20 天后还剩原来的几分之几？画出“剩余量—时间”的折线草图，在图上标出每个半衰期对应的点。能用“4 乘 5 等于 20，所以正好衰变完”来推理吗？错在哪里？
\item 用本章的三种衰变知识分析：为什么 $\alpha$ 放射源放在桌上时不必过度恐慌，而一旦把它研成粉末吸入体内就非常危险？从“剂量”和“射线种类”两个角度各说一条理由。
\item 画一幅“拔河图”：左侧画一个小核、右侧画一个很大的核，用箭头表示强相互作用（只连邻居）和电磁排斥（连向所有质子），解释为什么小核稳定、过大的核必然不稳定。图旁注明：这是人为的表示法，真实核里没有棍子和绳子。
\item 有人说：“核辐射太可怕了，我们应该生活在零辐射的环境里。”用“放射性、剂量、风险”三分法写一段反驳：零辐射在地球上是否可能？把天然本底降为零是否值得追求？
\item 一具考古出土的骨骼中，碳-14 含量只剩活体水平的八分之一。估算它的年代，并说明这个方法隐含的两个前提（提示：想想活体碳-14 水平从哪来、死亡后发生了什么）。
\end{sikaoti}
